% Source-derived chapter 06: Proof of Theorem~\ref{thm:E} and Theorem~\ref{thm:F}
% Original source: anticyclotomic-iwasawa-rational-elliptic-eisenstein
\chapter{Proof of Theorem~\ref{thm:E} and Theorem~\ref{thm:F}}
\label{anticyclotomic-iwasawa-rational-elliptic-eisenstein-chapter-06}
\source{anticyclotomic-iwasawa-rational-elliptic-eisenstein}

\begin{remark}[Source section]
\label{anticyclotomic-iwasawa-rational-elliptic-eisenstein-chapter-06-source}
\source{anticyclotomic-iwasawa-rational-elliptic-eisenstein}
\source{anticyclotomic-iwasawa-rational-elliptic-eisenstein}
This chapter reproduces the corresponding section of the retrieved
LaTeX source, including its definitions, statements, examples, and
proofs. It has not yet been translated into Lean.
\notready
\end{remark}

% The source section title was: Proof of Theorem~\ref{thm:E} and Theorem~\ref{thm:F}
\label{sec:EF}
\source{anticyclotomic-iwasawa-rational-elliptic-eisenstein}

\subsection{Preliminaries}

Here we collect the auxiliary results we shall use in the next  sections to deduce Theorems~\ref{thm:E} and~\ref{thm:F} in the introduction from our main result, Theorem~\ref{thm:thmA}.

\subsubsection{Anticyclotomic control theorem}

Assume that $p=v\bar{v}$ splits in $K$, and as in \cite[\S{2.2.3}]{jsw}, define the \emph{anticyclotomic Selmer group} of $W=E[p^\infty]$ by
\[
\rH^1_{\Fcal_{\rm ac}}(K,W)={\rm ker}\biggl\{
\rH^1(K^\Sigma/K,W)\rightarrow\prod_{w\in \Sigma}\rH^1(K_w,W)\times\frac{\rH^1(K_v,W)}{\rH^1(K_v,W)_{\rm div}}\times\rH^1(K_{\bar{v}},W)\biggr\},
\]
where $\rH^1(K_v,W)_{\rm div}\subset\rH^1(K_v,W)$ denotes the maximal divisible submodule and $\Sigma=\{w: \ w\vert N\}$. 

The following result is a special case of the ``anticyclotomic control theorem'' of \cite[\S{3}]{jsw}.


\begin{theorem}
\source{anticyclotomic-iwasawa-rational-elliptic-eisenstein}
\notready\label{controlthm}
\source{anticyclotomic-iwasawa-rational-elliptic-eisenstein}
Assume that
\begin{itemize}
%\item{} $E(K)[p]=0$ and 
\item $E(\bQ_p)[p]=0$,
\item{} ${\rm rank}_{\Z}E(K)=1$,
\item{} $\#\sha(E/K)[p^\infty]<\infty$.
\end{itemize}
Then $\mathfrak{X}_E$ is a torsion $\Lambda$-module, and letting $\Fcal_E\in\Lambda$ be a generator of ${\rm char}_\Lambda(\mathfrak{X}_E)$, we have
\[
\# \Z_p /\Fcal_E(0) = 
\# \sha(E/K)[p^\infty]\cdot\biggl(  \dfrac{ \# (\Z_p / ( \frac{1-a_p+p}{p}) \cdot \mathrm{log}_{\omega_E} P) }{ [E(K): \Z\cdot P]_p }\biggr)^2\cdot\prod_{w\mid N} c_w^{}(E/K)_p,
%\# \mathrm{H}^1_{\mathrm{ur}}(K_w, E[p^\infty])
\]	
where
\begin{itemize}
\item $P\in E(K)$ is any point of infinite order, 
\item $\log_{\omega_E}:E(K_v)_{/{\rm tors}}\rightarrow\bZ_p$ is the formal group logarithm associated to a N\'{e}ron differential $\omega_E$, 
\item $[E(K):\Z\cdot P]_p$ denotes the $p$-part of the index $[E(K):\Z\cdot P]$,
\item $c_w^{}(E/K)_p$ is the $p$-part of the Tamagawa number of $E/K_w$.
\end{itemize}
\end{theorem}

\begin{proof}
This follows from the combination of Theorem~3.3.1 and equation (3.5.d) in \cite[(3.5.d)]{jsw}, noting that the arguments in the proof of those results apply without change with the $G_K$-irreducibility of $E[p]$ assumed in \emph{loc.\,cit.} replaced by the weaker hypothesis that $E(K)[p]=0$, {which 
is implied by the hypothesis $E(\bQ_p)[p]=0$ since $p$ splits in $K$}. 
\end{proof}

 
\subsubsection{Gross--Zagier formulae}

Let $E/\bQ$ be an elliptic curve of conductor $N$, and fix a  parametrization
\[
\pi:X_0(N)\rightarrow E.
\]

Let $K$ be an imaginary quadratic field 
%of discriminant $D_K$ 
satisfying the Heegner hypothesis relative to $N$, and fix an integral ideal $\mathfrak{N}\subset\mathcal{O}_K$ with $\mathcal{O}_K/\mathfrak{N}=\Z/N\Z$. Let $x_1=[\mathbb{C}/\mathcal{O}_K\rightarrow\mathbb{C}/\mathfrak{N}^{-1}]\in X_0(N)$ 
be the Heegner point of conductor $1$ on $X_0(N)$, which is defined over the Hilbert class field $H=K[1]$ of $K$, and set
\[
P_K=\sum_{\sigma\in{\rm Gal}(H/K)}\pi(x_1)^\sigma\in E(K).
\]

Let $f\in S_2(\Gamma_0(N))$ be the newform associated with $E$, so that $L(f,s)=L(E,s)$, and consider the differential $\omega_f:=2\pi i f(\tau)d\tau$ on $X_0(N)$. Let also $\omega_E$ be a N\'{e}ron differential on $E$, and let $c_E\in\mathbb{Z}$ be the associated \emph{Manin constant}, so that $\pi^*(\omega_E)=c_E\cdot\omega_f$.
  
\begin{theorem}
\source{anticyclotomic-iwasawa-rational-elliptic-eisenstein}
\notready
\label{thmGZ} 
\source{anticyclotomic-iwasawa-rational-elliptic-eisenstein}
Under the above hypotheses, $L(E/K,1)=0$ and 
\begin{displaymath}
L'(E/K,1)=u_K^{-2}c_E^{-2}\cdot\sqrt{\vert D_K\vert}^{-1}\cdot\Vert \omega_E \Vert^2\cdot \hat{h}(P_K),
\end{displaymath}
where $u_K=\#(\mathcal{O}_K^\times/{\pm{1}})$, %$D_K<0$ is the discriminant of $K$, 
$\hat{h}(P_K)$ is the canonical height of $P_K$, and $\Vert\omega_E\Vert^2=\iint_{E(\C)}|\omega_E \wedge \bar{\omega}_E|$. 
\end{theorem}

\begin{proof}
This is \cite[Thm.~V.2.1]{grosszagier}.
\end{proof}

\begin{theorem}
\source{anticyclotomic-iwasawa-rational-elliptic-eisenstein}
\notready
\label{thmpadicGZ} 
\source{anticyclotomic-iwasawa-rational-elliptic-eisenstein}
Under the above hypotheses, let $p>2$ be a prime of good reduction for $E$ such that $p=v\bar{v}$ splits in $K$. Then
\[
\mathcal{L}_E(0)=c_E^{-2}\cdot\bigl(1-a_pp^{-1}+p^{-1}\bigr)^2\cdot{\rm log}_{\omega_E}(P_K)^2.
\]
where ${\rm log}_{\omega_E}:E(K_v)\rightarrow K_v$ is the formal group logarithm associated to $\omega_E$.
\end{theorem}

\begin{proof}
Let $J_0(N)$ be the Picard variety of $X_0(N)$, and set $\Delta_1=(x_1)-(\infty)\in J_0(N)(H)$. 
By \cite[Thm.~5.13]{BDP} specialized to the case $k=2$, $r=j=0$, and $\chi=\mathbf{N}_K^{-1}$, we have
\[
\mathcal{L}_E(0)=\bigl(1-a_pp^{-1}+p^{-1}\bigr)^2\cdot\biggl(\sum_{\sigma\in{\rm Gal}(H/K)}\log_{\omega_f}(\Delta_1^\sigma)\biggr)^2,
\]
where $\log_{\omega_f}:J_0(N)(H_v)\rightarrow H_v$ is the formal group logarithm associated to $\omega_f$. Since $\log_{\omega_f}(\Delta_1)=%c_E^{-1}\cdot\log_{\pi^*\omega_E}(\Delta_1)=
c_E^{-1}\cdot\log_{\omega_E}(\pi(\Delta_1))$, this yields the result.
\end{proof}



\subsubsection{A result of Greenberg--Vatsal}


\begin{theorem}
\source{anticyclotomic-iwasawa-rational-elliptic-eisenstein}
\notready
\label{thmGV}
\source{anticyclotomic-iwasawa-rational-elliptic-eisenstein}
Let $A/\Q$ be an elliptic curve, and let $p>2$ be a prime of good ordinary reduction for $A$. Assume that $A$ admits a cyclic $p$-isogeny with kernel $\Phi_A$, with the $G_\bQ$-action on $\Phi_A$ given by a character which is either ramified at $p$ and even, or unramified at $p$ and odd. If $L(A,1)\neq 0$ then
\begin{displaymath}
\ord_{p}\biggl(\frac{L(A, 1)}{\Omega_{A}}\biggr)=\ord_{p}\biggl(\frac{\# \sha(A / \Q)\cdot{\rm Tam}(A/\Q)}{\#(A(\bQ)_{\rm tors})^2}\biggr),
\end{displaymath}
where ${\rm Tam}(A/\bQ)=\prod_{\ell} c_\ell(A/\bQ)$  
is the product over the bad primes $\ell$ of $A$ of the Tamagawa numbers of $A/\bQ_\ell$.
\end{theorem}

\begin{proof}
By \cite{kolyvagin-mw-sha}, if $L(A,1)\neq 0$ then ${\rm rank}_\Z A(\Q)=0$ and $\#\sha(A/\Q)<\infty$; in particular, $\#{\rm Sel}_{p^\infty}(A/\bQ)=\#\sha(A/\bQ)[p^\infty]<\infty$. Letting $\Lambda_{\rm cyc}=\Z_p\llbracket{\rm Gal}(\Q_\infty/\Q)\rrbracket$ be the cyclotomic Iwasawa algebra, by \cite[Thm.~4.1]{greenberg-cetraro} we therefore have
\begin{equation}\label{eq:greenberg-4.1}
\source{anticyclotomic-iwasawa-rational-elliptic-eisenstein}
\#\Z_p/\mathcal{F}_A(0)=\frac{\#\bigl(\Z_p/(1-a_p(A)+p)^2\cdot\#\sha(A/\Q)\cdot{\rm Tam}(A/\bQ)\bigr)}{\#\bigl(\Z_p/(A(\Q)_{\rm tors})^2\bigr)},
\end{equation}
where $\mathcal{F}_A\in\Lambda_{\rm cyc}$ is a generator of the characteristic ideal of the dual Selmer group $Sel_{\Q_\infty}(T_pA,T_p^+A)^\vee$ in the notations of \cite[\S{3.6.1}]{skinner-urban}. Under the given assumptions, the cyclotomic main conjecture for $A$, i.e., the equality
\begin{equation}\label{eq:cyc-IMC}
\source{anticyclotomic-iwasawa-rational-elliptic-eisenstein}
(\mathcal{F}_A)=(\mathcal{L}_A)\subset\Lambda_{\rm cyc}
\end{equation}
where $\mathcal{L}_A$ is the $p$-adic $L$-function of Mazur--Swinnerton-Dyer, follows from the combination of \cite[Thm.~12.5]{kato-euler-systems} and\cite[Thm.~1.3]{greenvats}. By the interpolation property of $\mathcal{L}_A$, 
\begin{equation}\label{eq:interp-MSW}
\source{anticyclotomic-iwasawa-rational-elliptic-eisenstein}
\mathcal{L}_A(0)=\bigl(1-\alpha_p^{-1}\bigr)^2\cdot\frac{L(A,1)}{\Omega_A},
\end{equation}
where $\alpha_p\in\Z_p^\times$ is the unit root of $x^2-a_p(A)x+p$. Noting that ${\rm ord}_p(1-a_p(A)+p)={\rm ord}_p(1-\alpha_p^{-1})$, the result thus follows from the combination of (\ref{eq:greenberg-4.1}), (\ref{eq:cyc-IMC}), and (\ref{eq:interp-MSW}).
\end{proof}



\subsection{Proof of the \texorpdfstring{$p$}{p}-converse} 

The next result is Theorem~\ref{thm:E} in the introduction.
Note that a result for $r=0$ can be obtained from the cyclotomic main conjecture proved by combining \cite[Thm.~1.3]{greenvats} and Kato's divisibility in \cite{kato-euler-systems}. However, our assumptions are less restrictive, since in \cite{greenvats} the character $\phi$ is assumed to be ramified at $p$ and even or unramified at $p$ and odd.
\begin{theorem}
\source{anticyclotomic-iwasawa-rational-elliptic-eisenstein}
\notready\label{thm:thmB}
\source{anticyclotomic-iwasawa-rational-elliptic-eisenstein}
Assume that $E[p]^{ss}=\mathbb{F}_p(\phi)\oplus\mathbb{F}_p(\psi)$ with $\phi\vert_{G_p}\neq\mathds{1},\omega$. Let $r\in\{0,1\}$. Then 
\[
{\rm corank}_{\Z_p}{\rm Sel}_{p^\infty}(E/\Q) = r 
%\left.
%\begin{array}{r}
%{\rm rank}_\Z E(\bQ)=r\\[0.2em]
%\#\sha(E/\bQ)[p^\infty]<\infty
%\end{array}
%\right\}
\quad\Longrightarrow\quad
{\rm ord}_{s=1}L(E,s)=r,
\]
and so ${\rm rank}_\Z E(\Q)=r$ and $\#\sha(E/\Q)<\infty$.
\end{theorem}

\begin{proof}
The proof of this result is a consequence of Corollary~\ref{cor:PR} for suitable choices of a quadratic imaginary field $K$ depending on $r\in\{0,1\}$.

First we suppose ${\rm corank}_{\Z_p}{\rm Sel}_{p^\infty}(E/\Q) = 1$. It follows from \cite[Theorem 1.5]{monsky} that the root number $w(E/\Q)=-1$.  Choose an imaginary quadratic field $K$ of discriminant $D_K$ such that
\begin{itemize}
\item[(a)] $D_K<-4$ is odd,
\item[(b)] every prime $\ell$ dividing $N$ splits in $K$,
\item[(c)] $p$ splits in $K$, say $p=v\bar{v}$,
\item[(d)] $L(E^K,1)\neq 0$.
\end{itemize}
The existence of such $K$ (in fact, of an infinitude of them) is ensured by \cite[Thm.~B.1]{twist}, since (a), (b), and (c) impose only a finite number of congruence conditions on $D_K$, and any $K$ satisfying (b) is such that $E/K$ has root number $w(E/K)=w(E/\Q)w(E^K/\bQ)=-1$, and therefore $w(E^K/\bQ)=+1$. By work of Kolyvagin \cite{kolyvagin-mw-sha} (or alternatively, Kato \cite{kato-euler-systems}), the non-vanishing of $L(E^K,1)$ implies that ${\rm Sel}_{p^\infty}(E^K/\Q)$ is finite and therefore
\[
{\rm corank}_{\Z_p}{\rm Sel}_{p^\infty}(E/K) = 1,
\]
and hence $E$ satisfies \eqref{eq:intro-sel}.  In particular, all the hypotheses of Corollary~\ref{cor:PR} hold. 
As in the proof of Corollary~\ref{cor:howard}, the condition that ${\rm corank}_{\Z_p}{\rm Sel}_{p^\infty}(E/K) = 1$ easily implies that 
$\ord_{\mathfrak{P}_0}(X_{\rm tors}) = 0$, and so it then follows from Corollary \ref{cor:PR} that the image of $\kappa_1$ in 
$\rH^1(K,T_pE)$ is non-zero. This implies that the Heegner point $P_K\in E(K)$ is non-torsion and hence, by the Gross--Zagier formula that $\ord_{s=1}L(E/K,s) = 1$. Since $L(E/K,s)= L(E,s)L(E^K,s)$ and $\ord_{s=1}L(E^K,s) = 0$ by the choice of $K$,
it follows that $\ord_{s=1}L(E,s) = 1$.

We now assume ${\rm corank}_{\Z_p}{\rm Sel}_{p^\infty}(E/\Q) = 0$. The result of Monsky used above implies in this case that $w(E/\Q)=+1$. We now choose an imaginary quadratic field $K$ satisfying the same conditions (a), (b), (c) as above, in addition to the condition
\begin{itemize}
\item[(d')] $\ord_{s=1}L(E^K,1)=1$.
\end{itemize} 
The existence of infinitely many such $K$ follows from \cite[Thm.~B.2]{twist}, since any $K$ satisfying (b) is such that $w(E^K/\Q)=-1$. The Gross--Zagier--Kolyvagin theorem %\eqref{eq:GZK} 
implies that ${\rm corank}_{\Z_p}{\rm Sel}_{p^\infty}(E^K/\Q) = 1$ and therefore 
\[
{\rm corank}_{\Z_p}{\rm Sel}_{p^\infty}(E/K) = 1.
\]
Thus, as above, $E$ satisfies \eqref{eq:intro-sel}, and we can apply Corollary~\ref{cor:PR} and the Gross--Zagier formula to obtain $\ord_{s=1}L(E/K,s) = 1$, which implies by our choice of $K$ that $L(E,1)\neq 0$. 
\end{proof}


Since the hypotheses of Theorem \ref{thm:thmB} imply $E(\Q)[p]= 0$, we see that ${\rm Sel}_{p^\infty}(E/\Q)[p]={\rm Sel}_p(E/\Q)$, whence the following mod $p$ version of the theorem.
%it is easily seen that 
%${\rm Sel}_p(E/\Q)\simeq (\Z/p\Z)^r$ for $r\in\{0,1\}$ implies ${\rm Sel}_{p^\infty}(E/\Q) \simeq (\Q_p/\Z_p)^r$, that is, ${\rm corank}_{\Z_p} {\rm Sel}_{p^\infty}(E/\Q) = r$. This yields a mod $p$ version of the theorem:

\begin{cor}
\source{anticyclotomic-iwasawa-rational-elliptic-eisenstein}
\notready\label{cor:thmB} Suppose $E$ is as in Theorem \ref{thm:thmB} and $r\in\{0,1\}$. Then 
\source{anticyclotomic-iwasawa-rational-elliptic-eisenstein}
\[
\dim_{\mathbb{F}_p}{\rm Sel}_p(E/\Q) = r 
\quad\Longrightarrow\quad
{\rm ord}_{s=1}L(E,s)=r,
\]
and so ${\rm rank}_\Z E(\Q)=r$ and $\#\sha(E/\Q)<\infty$.
\end{cor}

For $p=3$, Corollary~\ref{cor:thmB} together with the work of Bhargava--Klagsbrun--Lemke Oliver--Shnidman \cite{BKLS} on the average $3$-Selmer rank in quadratic twist families, leads to the following result in the direction of Goldfeld's conjecture \cite{goldfeld}. 

\begin{cor}
\source{anticyclotomic-iwasawa-rational-elliptic-eisenstein}
\notready\label{cor:goldfeld}
\source{anticyclotomic-iwasawa-rational-elliptic-eisenstein}
Let $E$ be an elliptic curve over $\bQ$ with a rational $3$-isogeny. Then %for $r\in\{0,1\}$ 
a positive proportion of quadratic twists of $E$ have algebraic and analytic rank equal to $1$
and a positive proportion of quadratic twists of $E$ have algebraic and analytic rank equal to $0$.
\end{cor}

\begin{proof}
Denote by $\phi:G_\bQ\rightarrow\mathbb{F}_3^\times=\mu_2$ the character giving the Galois action on the kernel of a rational $3$-isogeny of $E$. As the condition $\phi\vert_{G_p}\neq\mathds{1},\omega$ can be arranged by a quadratic twist, combining \cite[Thm.~2.6]{BKLS} and Corollary~\ref{cor:thmB}, the result follows.
\end{proof}

\begin{remark}
\source{anticyclotomic-iwasawa-rational-elliptic-eisenstein}
\notready\label{rem:BKLS}
\source{anticyclotomic-iwasawa-rational-elliptic-eisenstein}
The qualitative result of Corollary~\ref{cor:goldfeld} was first obtained by Kriz--Li (see \cite[Thm.~1.5]{kriz-li}), but thanks to \cite{BKLS} (see esp. [\emph{op.\,cit.}, p.~2957]) our result {can lead} to better lower bounds on the proportion of rank~1 twists. In particular the proportions provided by \cite[Thm.~2.5]{BKLS} are the largest when the parity of the logarithmic Selmer ratios is equidistributed in quadratic families. The elliptic curve of smallest conductor over $\Q$ for which this happens is the elliptic curve having Cremona label $19a3$ given by the affine equation $y^2+y=x^3+x^2+x$. The explicit bounds of \cite{BKLS} and our result give that at least $41.6\%$ of its quadratic twists have analytic and algebraic rank equal to $1$ and at least $25\%$ have analytic and algebraic rank equal to $0$.
\end{remark}

\subsection{Proof of the \texorpdfstring{$p$}{p}-part of BSD formula} 

The following is Theorem~\ref{thm:F} in the introduction.

\begin{theorem}
\source{anticyclotomic-iwasawa-rational-elliptic-eisenstein}
\notready\label{thm:thmC}
\source{anticyclotomic-iwasawa-rational-elliptic-eisenstein}
Let $E/\bQ$ be an elliptic curve, and let $p>2$ be a prime of good ordinary reduction for $E$. Assume that $E$ admits a cyclic $p$-isogeny with kernel $C=\mathbb{F}_p(\phi)$, with $\phi:G_\bQ\rightarrow\mathbb{F}_p^\times$ such that
\begin{itemize}
\item $\phi\vert_{G_p}\neq\mathds{1},\omega$,
\item $\phi$ is either ramified at $p$ and odd, or unramified at $p$ and even.
\end{itemize}
If ${\rm ord}_{s=1}L(E,s)=1$, then 
\[
\ord_{p}\biggl(\frac{L'(E, 1)}{\operatorname{Reg}(E / \Q) \cdot \Omega_{E}}\biggr)=\ord_{p}\biggl(\# \sha(E / \Q) \prod_{\ell \nmid \infty} c_{\ell}(E / \Q)\biggr).
\]
In other words, the $p$-part of the Birch--Swinnerton-Dyer formula for $E$ holds.
\end{theorem}



\begin{proof}
Suppose ${\rm ord}_{s=1}L(E,s)=1$ and choose, as in the proof of Theorem~\ref{thm:thmB}, an imaginary quadratic field $K$ of discriminant $D_K$ such that
\begin{itemize}
\item[(a)] $D_K<-4$ is odd,
\item[(b)] every prime $\ell$ dividing $N$ splits in $K$,
\item[(c)] $p$ splits in $K$, say $p=v\bar{v}$,
\item[(d)] $L(E^K,1)\neq 0$.
\end{itemize}
Then ${\rm ord}_{s=1}L(E/K,s)=1$, which by Theorem~\ref{thmGZ} implies that the Heegner point $P_K\in E(K)$ has infinite order, and therefore ${\rm rank}_\Z E(K)=1$ and $\#\sha(E/K)<\infty$ by \cite{kolyvagin-mw-sha}. 
In particular, \eqref{eq:intro-sel} holds, and so all the hypotheses of Theorem \ref{thm:thmA} are satisfied. Thus there is a $p$-adic unit $u\in(\Z_p^{{\rm ur}})^\times$ for which
\begin{equation}\label{eq:thmA}
\source{anticyclotomic-iwasawa-rational-elliptic-eisenstein}
\mathcal{F}_E(0)=u\cdot\mathcal{L}_E(0),
\end{equation}
where $\mathcal{F}_E\in\Lambda$ is a generator of ${\rm char}_\Lambda(\mathfrak{X}_E)$. The hypotheses on $\phi$ imply that $E(K)[p]=0$, and so Theorem~\ref{controlthm} applies with $P=P_K$, which combined with Theorem~\ref{thmpadicGZ} and the relations $(\ref{eq:comparison})$ and $(\ref{eq:thmA})$ 
yields the equality
\begin{equation}\label{eq:thmA-bis}
\source{anticyclotomic-iwasawa-rational-elliptic-eisenstein}
{\rm ord}_p\bigl(\#\sha(E/K)\bigr)=2\;{\rm ord}_p\bigl(c_E^{-1}u_K^{-1}\cdot[E(K):\Z.P_K]\bigr)-\sum_{w\in S}{\rm ord}_p\bigl(c_w(E/K)\bigr),
\end{equation}
 
On the other hand, the Gross--Zagier formula of Theorem~\ref{thmGZ} can be rewritten (see \cite[p.~312]{grosszagier}) as
\[
L'(E/K,1)=2^tc_E^{-2}u_K^{-2}\cdot\hat{h}(P_K)\cdot
\Omega_E\cdot\Omega_{E^K},
\]
where the power of $2$ is given by the number of connected components $[E(\mathbb{R}):E(\mathbb{R})^0]$.
This, together with the relations $L(E/K,s)=L(E,s)\cdot L(E^K,s)$ and
\[
\hat{h}(P_K)=[E(K):\Z\cdot P_K]^2\cdot{\rm Reg}_{}(E/K)=[E(K):\Z\cdot P_K]^2\cdot{\rm Reg}_{}(E/\Q),
\]
using that ${\rm rank}_\Z E^K(\Q)=0$ for the last equality, amounts to the formula
\begin{equation}\label{eq:GZ}
\source{anticyclotomic-iwasawa-rational-elliptic-eisenstein}
\frac{L'(E,1)}{{\rm Reg}(E/\bQ)\cdot\Omega_E}\cdot\frac{L(E^K,1)}{\Omega_{E^K}}=2^tc_E^{-2}u_K^{-2}\cdot[E(K):\Z\cdot P_K]^2.
\end{equation}
Note that $u_K=1$, since $D_K<-4$. Since $\sha(E/K)[p^\infty] \simeq\sha(E/\bQ)[p^\infty]\oplus\sha(E^K/\bQ)[p^\infty]$ as $p$ is odd, and   
\[
\sum_{w\vert\ell}{\ord}_p(c_w(E/K))={\rm ord}_p(c_\ell(E/\bQ))+{\rm ord}_p(c_\ell(E^K/\Q))
\]
for any prime $\ell$ (see \cite[Cor.~9.2]{skinner-zhang}), combining $(\ref{eq:thmA-bis})$ and (\ref{eq:GZ}) we arrive at
\begin{equation}\label{eq:thmA-bisbis}
\source{anticyclotomic-iwasawa-rational-elliptic-eisenstein}
\begin{aligned}
{\rm ord}_p\biggl(\frac{L'(E,1)}{{\rm Reg}(E/\Q)\cdot\Omega_E\cdot\prod_\ell c_\ell(E/\bQ)}\biggr)&-{\rm ord}_p(\#\sha(E/\Q))\\
&=
{\rm ord}_p\biggl(\frac{L(E^K,1)}{\Omega_{E^K}\cdot\prod_\ell c_\ell(E^K/\Q)}\biggr)-{\rm ord}_p\bigl(\#\sha(E^K/\Q)\bigr).
\end{aligned}
\end{equation}

Finally, by our hypotheses on $\phi$ the curve $E^K$ satisfies the hypotheses of Theorem~\ref{thmGV}, and hence the right-hand side of $(\ref{eq:thmA-bisbis})$ vanishes, concluding the proof of Theorem~\ref{thm:thmC}.
\end{proof} 



\bibliographystyle{alpha}
\bibliography{references}
